\documentclass[dvipdfmx, uplatex]{jsarticle}
\usepackage[dvipdfmx]{graphicx}
\usepackage{amsmath}
\usepackage{amssymb}
\usepackage{amsfonts}
\usepackage{amsthm}
\usepackage{mathrsfs}
\usepackage{bm}
\usepackage{braket}

\usepackage{silence}
\WarningFilter*{hyperref}{Token not allowed in a PDF string}
\usepackage[dvipdfmx]{hyperref}
\usepackage{pxjahyper}
\allowdisplaybreaks[4]

\usepackage[left=25truemm,top=20truemm,bottom=20truemm,right=25truemm]{geometry}

\newcommand{\bk}{\bm{k}}
\newcommand{\bx}{\bm{x}}

\theoremstyle{definition}
\newtheorem{prob}{問題}

\title{演習：正方格子上のランダムウォークと拡散極限\\
\large ―― 離散確率過程・特性関数・diffusive scaling ――}
\author{}
\date{}

\begin{document}
\maketitle
\vspace{-6zh}

\begin{abstract}
無限 2 次元正方格子上を歩く粒子の離散確率過程（ランダムウォーク）を題材に，
master 方程式の立式，特性関数（Fourier 変換）による解，平均二乗変位，そして格子間隔 $a$・
時間刻み $\tau$ を $0$ へ送る連続極限を扱う。鍵は，$a$ だけ，あるいは $\tau$ だけを $0$ にする
素朴な極限が\emph{自明}になってしまうこと，非自明な拡散過程を得るには $\tau$ を $a$ の関数として
両者を同時に $0$ へ送る\textbf{スケーリング（diffusive scaling）}が必要で，そこから拡散係数 $D$ が
現れること，を自分で見つけることにある。さらに，格子点間の自然な距離が\textbf{マンハッタン距離}
（$L^1$）であるにもかかわらず，連続極限の等確率面が\textbf{円}（$L^2$，等方的）になることを，
経路数（エントロピー）と中心極限定理＝Stirling の公式から導く――格子の離散対称性（$90^\circ$ 回転
$C_4$）から連続の回転対称性 $O(2)$ が\emph{創発}する様子を見る（この異方的スケーリング $z=2$ の構造は
続編 \texttt{exercise\_diffusion\_extra\_dimension} のスケーリング次元の議論に引き継がれる）。
最後に，得られる偏微分方程式の\emph{微分の階数}が何によって決まるかを考察する。
\textbf{各問は表式・数値を答えさせる形にしてあり，本文中には
原則として解答を与えていない}（正しく計算できていれば表式が一致するはずである）。完全な解答は
別ファイル \texttt{exercise\_lattice\_random\_walk\_solutions.tex} にある。
\end{abstract}

\tableofcontents

%======================================================================
\section{設定}
%======================================================================
格子間隔 $a$ の無限 2 次元正方格子を考える。格子点 $(i,j)$（$i,j\in\mathbb{Z}$）は
位置 $\bx=(ia,\,ja)$ にある。時間は離散的で，1 ステップを $\tau$ とする（時刻 $t=n\tau$）。
粒子は時刻 $0$ に原点にあり，各ステップごとに独立に
\begin{itemize}
\item 確率 $p$ で隣接する 4 格子点のいずれか 1 つへ等確率に（各 $\tfrac{p}{4}$ で）移動し，
\item 確率 $1-p$ でその場に留まる
\end{itemize}
（$0<p\le1$）。時刻 $n\tau$（$n$ ステップ後）に格子点 $(i,j)$ に粒子が存在する確率を
$P_n(i,j)$ と書く。

%======================================================================
\section{問題}
%======================================================================

\begin{prob}\label{prob:rw}
\begin{enumerate}
\item \textbf{初期条件。}粒子が原点に局在する初期条件 $P_0(i,j)$ を式で書け
（Kronecker の $\delta$ を用いよ）。これが表す状況を一言で述べよ。
\item \textbf{master 方程式。}1 ステップの規則から，$P_{n+1}(i,j)$ を $P_n$ の
（自分自身と近傍 4 点の）値で表す漸化式を書け。また確率の総和 $\sum_{i,j}P_n(i,j)$ が
$n$ によらず一定であることを確かめ，その値を述べよ。
\item \textbf{特性関数（Fourier 変換）。}特性関数を
$\displaystyle\hat P_n(\bk):=\sum_{i,j}P_n(i,j)\,e^{i\bk\cdot\bx}
=\sum_{i,j}P_n(i,j)\,e^{i(k_x ia+k_y ja)}$
で定める。(2) の漸化式から，これが 1 段の関係 $\hat P_{n+1}(\bk)=\lambda(\bk)\,\hat P_n(\bk)$ を
満たすことを示し，\textbf{$\lambda(\bk)$ を $\cos$ を用いて表せ}。さらに初期条件から $\hat P_0(\bk)$ を求め，
\textbf{$\hat P_n(\bk)$ を $\lambda(\bk)$ と $n$ で表せ}。
\item \textbf{平均二乗変位。}1 ステップあたりの変位 $\Delta\bx$ について $\langle\Delta x^2\rangle$ を
$p,a$ で求めよ。これを用いて（各ステップが独立であることから），$n$ ステップ後の平均二乗変位
$\langle x^2+y^2\rangle_n$ を $n,p,a$ で表せ。
\item \textbf{素朴な極限は自明であること。}$t=n\tau$（すなわち $n=t/\tau$）を固定して連続極限を考える。
(4) で求めた平均二乗変位 $\langle x^2+y^2\rangle$（物理的な長さの単位で）を用いて，次の 2 つの素朴な極限で
分布の空間的な広がりがどうなるかを調べ，いずれも自明（退化した）極限になることを述べよ：
\textbf{(i)} $\tau$ を固定したまま $a\to0$；\textbf{(ii)} $a$ を固定したまま $\tau\to0$。
\item \textbf{diffusive scaling と拡散係数。}(5) を踏まえ，連続極限 $a,\tau\to0$ で分布の広がりを
\emph{有限かつ非ゼロ}に保つには，$\tau$ を $a$ のどのような関数として（$a$ と $\tau$ をどんな関係に保って）
$0$ へ送ればよいかを求めよ（diffusive scaling）。その極限で，平均二乗変位が
$\langle x^2+y^2\rangle=(\text{定数})\times t$ となるときの比例係数から拡散係数 $D$ を定義し，
\textbf{$D$ を $p,a,\tau$ で表せ}。
\item \textbf{連続極限の方程式と解。}
\begin{enumerate}
\item master 方程式を $P_{n+1}(i,j)-P_n(i,j)$ の形に書き，$a,\tau$ が小さいとして
$P(x,y,t)$ の Taylor 展開で両辺を評価せよ。(6) の diffusive scaling のもとで $a,\tau\to0$ とし，
$P(x,y,t)$ が満たす偏微分方程式を導け。
\item (3) の $\hat P_n(\bk)$ の diffusive 極限 $\displaystyle\lim\hat P_n(\bk)$ を $D,t,\bk$ で求めよ
（$\lambda(\bk)$ の小 $\bk$ 展開を使う）。これを逆 Fourier 変換して，連続極限の確率密度 $P(x,y,t)$ を求めよ。
\item 得た $P(x,y,t)$ が (a) の偏微分方程式を満たすことを確かめよ。さらに連続極限の平均二乗変位
$\langle x^2+y^2\rangle$ を $D,t$ で求め，(4)(6) の結果と整合することを確かめよ。
\end{enumerate}
\item \textbf{マンハッタン距離と等確率面――等方性の創発。}格子上では，2 点を結ぶ経路は
軸に沿ってしか進めないため，最短経路の長さは\textbf{マンハッタン距離}（$L^1$ 距離）で測られる。
それにもかかわらず連続極限の等確率面が\emph{円}になることを，経路数（多重度＝エントロピー）と
Stirling の公式から確かめる。
\begin{enumerate}
\item 原点から格子点 $(i,j)$ へ到達する\emph{最小ステップ数}を $i,j$ で表せ（``stay'' は使わないとする）。
もし確率が最短経路だけで支配される（長い経路が指数的に抑制される）状況だったとしたら，
等確率面 $P=$ 一定はどんな図形（どの距離 $=$ 一定の面）になるか，$x,y$ の関係式で表せ。
\item まず 1 次元で考える。$N$ ステップ中 $k$ 回が $+a$（残りが $-a$）の歩みが終点 $x=(2k-N)a$ に
与える経路数（多重度）$\Omega$ を二項係数で表し，Stirling の公式
$\ln m!\simeq m\ln m-m+\tfrac12\ln(2\pi m)$ を用いて $\ln\Omega$ を（$k=N/2$ まわりで $x$ の 2 次まで）
展開せよ。ピークからの\emph{エントロピー欠損} $\Delta S:=\ln\Omega_{\max}-\ln\Omega$ を $x,N,a$ で表し，
対応する Gauss 分布の分散 $\sigma_x^2$ を読み取れ。
\item 2 次元にもどる。(7)(b) で示した連続極限の特性関数の振る舞いから，$P(x,y,t)$ が $x$ 部分と
$y$ 部分の積（独立な Gauss の積）になることを使い，等確率面 $P(x,y,t)=$ 一定を $x,y$ の関係式で表せ。
それはどんな図形か。また分散 $\sigma_x^2=\sigma_y^2$ を $D,t$ で表せ。
\item \textbf{考察。}なぜ等確率面はマンハッタン距離の菱形（$L^1$ 球）ではなく円（$L^2$ 球）に
なるのか。確率を支配するのが最短経路（``エネルギー''）ではなく経路数（``エントロピー''）であること，
中心極限定理によりその対数が等方的な 2 次形式になること，そして $x,y$ 方向の 1 ステップ分散が
等しく相関がないこと，の 3 点に触れて説明せよ。格子の離散対称性（$90^\circ$ 回転 $C_4$）から
連続の回転対称性 $O(2)$ が\emph{創発}的に拡大する，という言葉でまとめよ。
\end{enumerate}
\item \textbf{考察：微分の階数。}(7)(a) で得た連続極限の方程式は，時間について 1 階・空間について
2 階の微分方程式であった。次を考察せよ。
\begin{enumerate}
\item 空間微分が 2 階であって 1 階でないのはなぜか。本問の歩みの対称性（各向きへの移動確率）と
1 ステップの平均変位 $\langle\Delta\bx\rangle$ に関連づけて説明せよ。
\item 「時間 1 階・空間 2 階」という微分の階数が，diffusive scaling（(6) で求めた $\tau$ と $a$ の関係）と
どのように結びつくかを述べよ。
\item もし歩みに偏り（drift）があって 1 ステップの平均変位が $\langle\Delta x\rangle\neq0$（$a$ のオーダー）で
あったなら，連続極限で現れる主要項の空間微分の階数はいくつになり，そのとき $\tau$ を $a$ の何乗で
スケールさせればよいか，どんな型の方程式（拡散型／移流型）になるかを論ぜよ。
\end{enumerate}
\end{enumerate}
\end{prob}

\bigskip
{\sloppy
\noindent\textbf{解答について：}各問の完全な解答（求めるべき表式・数値を含む）は別ファイル
\texttt{exercise\_lattice\_random\_walk\_solutions.tex} にある。
\par}

\end{document}
